%!TeX root = ../master.tex
\documentclass[../master.tex]{subfiles}

% NOTE: If compiling this subfile completely by itself, the following 
% fallback packages will automatically initialize.
\makeatletter
\@ifundefined{lstlisting}{
    \usepackage{listings}
    \usepackage{xcolor}
    \usepackage{booktabs}
    \usepackage{amsmath}
}{}
\makeatother

\begin{document}

\chapter{Asynchronous Outbound Emergency Reporting Subsystem}
\label{ch:emergency_reporting}

\section{Subsystem Architectural Overview}
To ensure automated facility-level life-safety response during critical environmental breaches, an **Asynchronous Outbound Telephony and Emergency Reporting Subsystem** was integrated into the \textit{Environment-Safety-Monitor} firmware architecture. 

When environmental metrics breach predefined hazard levels, the system automatically computes physical gas displacement volumes and passes an automated Text-to-Speech (TTS) payload down a decoupled, non-blocking execution path. This data payload is then securely delivered over a public switched telephone network (PSTN) via cloud telephony gateways (e.g., Twilio) or localized Session Initiation Protocol (SIP/VoIP) trunks, all without interfering with the server's strict deterministic execution parameters.

\section{Mathematical Modeling: Asphyxiant Gas Displacement}
Because local electrochemical or optical sensors record environmental parameters as relative percentage concentrations ($O_2\%$) rather than direct volumes, the software must calculate the total volume of an introduced asphyxiant (e.g., Nitrogen, $N_2$) dynamically. 

Assuming a standard nominal baseline atmospheric composition of approximately $21.0\%$ Oxygen ($O_{2,\text{baseline}}$), a relative drop in measured oxygen indicates a proportional volumetric displacement by the escaping target gas. Let $V_{\text{room}}$ represent the fixed, enclosed three-dimensional volume of the monitored containment space in cubic meters ($m^3$), and $O_{2,\text{measured}}$ represent the real-time calculated oxygen concentration. The fraction of oxygen displaced, $\Delta\Phi_{O_2}$, is mathematically modeled as:

\begin{equation}
    \Delta\Phi_{O_2} = 1.0 - \left( \frac{O_{2,\text{measured}}}{O_{2,\text{baseline}}} \right)
\end{equation}

By scaling this fractional displacement against the total volumetric bounds of the physical room architecture, the runtime system derives the absolute volume of released nitrogen, $V_{N_2,\text{released}}$:

\begin{equation}
    V_{N_2,\text{released}} = \max\left(0.0, \, V_{\text{room}} \times \left( 1.0 - \frac{O_{2,\text{measured}}}{21.0} \right)\right)
\end{equation}

This calculated metric is passed directly to the emergency reporting module to provide public safety dispatchers (PSAPs) with an accurate assessment of the hazard scale before emergency units arrive on site.

\section{Hardware Interface Loop Metrics}
To guarantee system survival under electrical fault conditions, the physical wiring bridge between the monitoring server and the Advanced Axis AX panel utilizes a supervised three-state loop interface via the 55000-820ADV module. The hardware metrics, calculated line impedances, and expected system state transitions are summarized in Table~\ref{tab:hardware_loop_specs}.

\begin{table}[htbp]
    \centering
    \caption{55000-820ADV Supervised Loop Impedance Metrics}
    \label{tab:hardware_loop_specs}
    \begin{tabular}{llll}
        \toprule
        \textbf{Physical State} & \textbf{Relay Position} & \textbf{Target Loop Impedance} & \textbf{System Action Response} \\
        \midrule
        Standby / Nominal & Normally Closed (N.C.) & $\approx 6.67\text{ k}\Omega$ ($10\text{k}\parallel20\text{k}$) & Normal Baseline Log Monitoring \\
        Emergency Breach & Open Circuit Trigger & $\approx 20.00\text{ k}\Omega$ ($20\text{k}$ EOL Only) & Trip Relay + Outbound VoIP Call \\
        Field Line Shear & Cut / Disconnected Wire & $\infty\ \Omega$ (Open Circuit Marker) & Master FACP Yellow Trouble LED \\
        Line Pinch Fault & Pinched / Grounded Wire & $\le 2.00\ \Omega$ (Direct Short Circuit) & Reject Trigger, Flag Error State \\
        \bottomrule
    \end{tabular}
\end{table}

\section{Asynchronous Core Multi-Threading \& Cooldown Latching}
Executing network socket bindings, DNS routing calls, or streaming media payloads directly within the primary safety loop introduces non-deterministic delays. This can easily exceed the tight $100\text{ ms}$ processing window managed on Core~0. To prevent thread blocking and shield the primary safety engine from latency spikes, the subsystem implements a **Lock-Free Single-Producer Single-Consumer (SPSC) Ring Buffer** paired with an independent background execution task worker.

\subsection{Asynchronous Cooldown Timer Control Architecture}
To prevent line saturation, accidental network flooding, or repetitive outbound dialing loops during a sustained emergency event, the processing logic on Core~1 utilizes an **Atomic Cooldown Latch**. 

Let $t_{\text{current}}$ represent the current high-resolution system clock timestamp, and $t_{\text{last\_call}}$ represent an atomic variable tracking the exact time point of the most recent outbound call dispatch. The configuration enforces a rigid cooldown duration window, $T_{\text{cooldown}}$ (e.g., $300\text{ seconds}$). An outbound emergency notification is cleared for execution if and only if:

\begin{equation}
    (t_{\text{current}} - t_{\text{last\_call}}) \ge T_{\text{cooldown}} \quad \land \quad \text{Worker}_{\text{Active}} = \text{False}
\end{equation}

\subsection{Summary of System Capabilities}
By combining physical gas laws with asynchronous software patterns, the tracking software achieves several key goals:
\begin{itemize}
    \item \textbf{Preserved Loop Determinism:} Core~0 functions with near-zero software jitter ($<1\text{ ms}$), completely insulated from network delivery handshakes.
    \item \textbf{Actionable Telemetry Dispatch:} Dispatched audio alert streams verbally transmit exact physical metrics, providing responders with real-time situational awareness ($O_2\%$ concentration and calculated $N_2$ gas release volumes).
    \item \textbf{Infrastructure Protection:} The atomic state gating and timing windows protect communication channels from resource starvation, keeping phone lines accessible for emergency personnel.
\end{itemize}

\nocite{nfpa72_2022, twilio_compliance_2024}
\bibliographystyle{ieeetr}
\bibliography{references}

\end{document}
